\chapter*{Introduction générale}
\addcontentsline{toc}{chapter}{Introduction générale}
Une hausse des prix modifie la dépense nécessaire pour acheter une quantité donnée de biens et services. Son effet sur un budget dépend de la composition des achats. Lorsque l'alimentation occupe une part importante des dépenses, une poussée des prix alimentaires peut peser davantage sur l'équilibre budgétaire. Les données marocaines permettent d'étudier ce mécanisme à deux niveaux. L'IPC renseigne sur les prix suivis par le Haut-Commissariat au Plan (HCP). L'Enquête nationale sur le niveau de vie des ménages (ENNVM) décrit la répartition des dépenses de profils agrégés \cite{hcpIPC,hcpENNVM}.

L'exposition budgétaire désigne ici la variation du coût d'un sous-panier de composition fixe lorsque les prix changent. L'IPC renseigne sur les variations de prix et l'ENNVM 2022--2023 fournit les pondérations à un moment donné. Leur rapprochement répond à une question de simulation : quelle dépense supplémentaire faudrait-il engager pour conserver les mêmes quantités ? Observer le comportement des consommateurs demanderait de suivre leurs achats dans le temps pour mesurer, par exemple, des substitutions, des réductions de quantités ou des renoncements.

\section*{Problématique et objectifs}
Comment les différences de structure des dépenses modifient-elles l'exposition aux hausses de prix des profils nationaux, urbains et ruraux et des différents quintiles de dépenses au Maroc ? Le rapport décrit d'abord la dynamique de l'IPC général et de ses divisions. Il rapproche ensuite les postes de l'ENNVM des séries IPC disponibles, construit des indices de coût à poids fixes, puis distingue l'évolution aux prix historiques de deux chocs hypothétiques. La comparaison des deux périmètres permet d'apprécier dans quelle mesure les résultats dépendent de la couverture et des rapprochements de catégories.

\section*{Hypothèses de travail}
\begin{enumerate}
\item Une hausse uniforme de tous les prix retenus produit le même taux de variation pour chaque sous-panier. Le supplément imputable au budget total varie avec la part couverte.
\item Un choc alimentaire supplémentaire accroît particulièrement la dépense nécessaire pour les profils dont la part alimentaire initiale est élevée. Le résultat suppose des quantités fixes et un même budget initial dans l'illustration en dirhams.
\item L'élargissement du panier peut modifier le taux calculé et sa contribution au budget initial, puisque les nouveaux postes ont des trajectoires de prix propres et modifient les poids relatifs.
\end{enumerate}
Ces hypothèses décrivent des mécanismes comptables et économiques. Les données permettent de vérifier leur traduction dans les simulations.
